(1) 真空中半径 \(a\) 的透明水滴折射率 \(1<n<2\)，光线入射角 \(i\) 相对于法线，球内折射角 \(r\)。考虑入射折射、一次内反射、出射折射的主彩虹光线，用 Snell 定律求总偏向角 \(D(i)\)。
(2) 求产生主彩虹的最小偏向驻点 \(i_*\)，证明 \(\sin^2i_*=(4-n^2)/3\)；计算二阶导数并证明是非退化折叠焦散。取 \(n=4/3\) 给出主彩虹相对反太阳点的角半径。
(3) 假设入射平行光束在冲量参数 \(b=a\sin i\) 上均匀照明，计算几何光学微分截面 \(d\sigma/d\Omega\) 在 \(D\downarrow D_*\) 时由驻点两侧光线贡献的最强奇性及其系数，检查面积量纲。
(4) 将波动场在折叠焦散附近局部展开为规范积分 \(\int_{-\infty}^{\infty}\exp[ik(\xi t+t^3/3)]dt\)（\(k>0\)，\(\xi\) 为带符号控制参数且已归一化），以 Airy 函数表示其振幅，说明为何消除几何光学发散；讨论衍射标度与透明液滴的吸收、偏振和有限尺寸局限。
